% Related Work — tightened to ~0.5 page for the 8-page workshop budget.
% Four paragraphs, one per positioning beat: offline RL for portfolios,
% GRPO and the discrete-action assumption, exogeneity-based critic
% removal, O2O conservatism scheduling. li2024grpoexp placeholder
% removed; surrounding sentence rewritten to not need it.

\section{Related Work}
\label{sec:related}

\paragraph{Offline RL for portfolio allocation.} Continuous-allocation
RL on financial data \citep{jiang2017deep,liu2020finrl,moody2001learning,
deng2017deep} has increasingly imported modern offline-RL methods
trained on behavior-policy-mixture buffers: BC \citep{pomerleau1991efficient},
BCQ \citep{fujimoto2019off}, CQL \citep{kumar2020conservative},
TD3+BC \citep{fujimoto2021minimalist}, IQL \citep{kostrikov2022offline},
AWAC \citep{nair2020awac}, and sequence-model variants
\citep{chen2021decision,an2021uncertainty}. Our milestone showed that
this stack collapses smoothly under realistic transaction-cost friction
on the 8-ETF universe; we re-establish that collapse across a broader
algorithm grid (Sec.~\ref{sec:experiments:phase2a}) and ask whether a
critic-free, friction-aware on-policy method side-steps the pathology
rather than mitigating it.

\paragraph{GRPO and the discrete-action assumption.} Group Relative
Policy Optimization \citep{shao2024deepseekmath} replaces PPO's learned
value baseline with a within-group statistic computed across $G$
sampled completions of the same prompt. Subsequent work has applied
GRPO almost exclusively to LLM post-training pipelines
\citep{yu2024rlhf}, where the original formulation's two structural
assumptions hold by construction: actions are discrete
tokens with a finite alphabet, and ``state transitions'' are
deterministic sequence-extension. Our continuous-action portfolio
setting violates the first assumption (the policy emits a Dirichlet
density over the simplex) and trivially satisfies a stronger version
of the second (state transitions are independent of the agent's
action; Eq.~\ref{eq:exogeneity}). The cost we pay for this transfer is
the exogeneity precondition --- our setting satisfies it and most
other RL settings do not. We treat that asymmetry as a design
constraint, not a generality claim, and provide unit tests that make
the constraint legible to anyone porting the method
(\texttt{tests/test\_exogeneity.py}).

\paragraph{Critic removal under exogenous transitions.} The exogeneity
argument we use to retire the value baseline (Sec.~\ref{sec:method:grpo:exogeneity})
is closely related to direct policy-gradient estimators in
contextual-bandit and counterfactual-policy-gradient settings
\citep{dudik2014doubly,joachims2018deep,foerster2018counterfactual}:
when state transitions do not respond to the agent's action, the
bootstrapped continuation cancels in within-state action comparisons.
The financial-portfolio variant is implicit in early policy-gradient
trading work \citep{moody2001learning}, but the connection to
group-relative advantage estimation as a critic-removal mechanism has
not, to our knowledge, been formalized.

\paragraph{O2O conservatism scheduling.} Offline-to-online fine-tuning
has been studied in robotics
\citep{nair2020awac,zheng2023adaptive,zhao2022adaptive}, where the
typical solution is a fixed schedule on the conservatism coefficient
(linear or cosine decay over a known online budget). Our adaptive
schedule (Sec.~\ref{sec:method:o2o}) closes the loop by gating
pessimism on a per-step regime-KL signal. The Geodesic-CQL component
\citep[Fisher-Rao penalty;][]{ma2022geodesic} contributes the offline
phase; we do not modify it. In contrast to prior O2O work, which frames
the schedule's job as governing online responsiveness to distribution
shift, our Phase~2C results suggest its load-bearing function on this
universe is convergence-basin selection during fine-tune --- a
mechanism revision we unpack in Sec.~\ref{sec:discussion}.
